\documentclass[11pt,conference]{IEEEtran}
\usepackage{amsmath,amssymb,amsthm}
\usepackage{mathtools}
\usepackage{algorithm}
\usepackage{algorithmic}
\usepackage{listings}
\usepackage{hyperref}
\usepackage{glossaries}
\usepackage{xcolor}
\usepackage{tikz}
\usetikzlibrary{automata,positioning,arrows}

% Code listing configuration
\lstset{
    language=C,
    basicstyle=\footnotesize\ttfamily,
    keywordstyle=\color{blue}\bfseries,
    commentstyle=\color{gray}\itshape,
    stringstyle=\color{red},
    numbers=left,
    numberstyle=\tiny\color{gray},
    frame=single,
    breaklines=true,
    captionpos=b
}

% Theorem environments
\theoremstyle{definition}
\newtheorem{definition}{Definition}
\newtheorem{theorem}{Theorem}
\newtheorem{lemma}{Lemma}
\newtheorem{proposition}{Proposition}

% Mathematical notation
\newcommand{\actors}{\mathcal{A}}
\newcommand{\states}{\mathcal{S}}
\newcommand{\real}{\mathbb{R}}
\newcommand{\nat}{\mathbb{N}}

% Make glossary
\makeglossaries

% Verb-Noun Glossary Entries
\newglossaryentry{verifytrace}{
    name={verify-trace-$\Phi$},
    description={Validate cryptographic integrity of state transition history with epistemic signature $\Phi$}
}

\newglossaryentry{emitrollback}{
    name={emit-rollback-$\Phi$},
    description={Generate rollback event with epistemic signature for state recovery}
}

\newglossaryentry{audittransition}{
    name={audit-transition-$\Phi$},
    description={Inspect state lifecycle compliance with confidence metrics}
}

\newglossaryentry{commitstate}{
    name={commit-state-$\Phi$},
    description={Finalize state persistence to DIRAM with receipt generation}
}

\newglossaryentry{blockprogress}{
    name={block-progress-$\Phi$},
    description={Halt execution pending upstream epistemic resolution}
}

\newglossaryentry{validateconfidence}{
    name={validate-confidence-$\Phi$},
    description={Assess proof\_confidence against 95.4\% threshold}
}

\newglossaryentry{computeratio}{
    name={compute-ratio-$\Phi$},
    description={Calculate success:failure metrics for cascade detection}
}

\newglossaryentry{tracedependency}{
    name={trace-dependency-$\Phi$},
    description={Map hierarchical state relationships for rollback scope}
}

\newglossaryentry{enforcethreshold}{
    name={enforce-threshold-$\Phi$},
    description={Apply epistemic boundaries to state transitions}
}

\newglossaryentry{allocatememory}{
    name={allocate-memory-$\Phi$},
    description={Create cryptographically governed memory allocation}
}

\newglossaryentry{generatereceipt}{
    name={generate-receipt-$\Phi$},
    description={Produce SHA-256 trace for forensic accountability}
}

\newglossaryentry{appendtrace}{
    name={append-trace-$\Phi$},
    description={Add state transition to immutable DIRAM log}
}

\newglossaryentry{anchorhardware}{
    name={anchor-hardware-$\Phi$},
    description={Bind epistemic state to physical memory substrate}
}

\newglossaryentry{initiatecascade}{
    name={initiate-cascade-$\Phi$},
    description={Trigger strategic rollback sequence}
}

\newglossaryentry{propagaterollback}{
    name={propagate-rollback-$\Phi$},
    description={Extend rollback to dependent state hierarchy}
}

\newglossaryentry{memoizedelta}{
    name={memoize-delta-$\Phi$},
    description={Store confidence degradation for future reference}
}

\newglossaryentry{recoverstate}{
    name={recover-state-$\Phi$},
    description={Restore system to pre-failure configuration}
}

\title{Hierarchical Actor-Orchestrated State Management with DIRAM-Backed Epistemic Validation}

\author{\IEEEauthorblockN{OBINexus Computing - Aegis Framework Division}
\IEEEauthorblockA{Technical Specification for Actor Sub-ConOps Architecture\\
Document Classification: Production Infrastructure\\
Compliance: NASA-STD-8739.8, AEGIS-PROOF-1.2}}

\begin{document}
\maketitle

\begin{abstract}
We present the hierarchical state resolution model for Actor-orchestrated systems, extending the OBIAI Actor class through sub-conceptual task decomposition with DIRAM-backed memory governance. Each EA Actor autonomously manages task lifecycles using a TO-DO $\rightarrow$ DOING $\rightarrow$ DONE progression model, maintaining epistemic validation at 95.4\% confidence threshold. The system implements strategic rollback cascades when success:failure ratios fall below 1:2, ensuring self-correcting behavior through cryptographically traced state transitions. This architecture represents deployed production infrastructure, not theoretical design, providing forensic-level accountability through SHA-256 receipt logs and verb-noun conceptual modeling aligned with the Actor class tuple $\alpha = (S, \mathcal{C}, \Phi, \Psi, \epsilon)$.
\end{abstract}

\section{Introduction}

The hierarchical state resolution model extends the Actor class defined in the OBIAI framework through systematic sub-conceptual decomposition. Building upon the categorical foundation where Actors navigate infinite-dimensional semantic manifolds, we implement a production-ready state management system that maintains epistemic discipline while enabling autonomous task orchestration.

\begin{definition}[Actor Class Extension]
Given an Actor $\alpha = (S, \mathcal{C}, \Phi, \Psi, \epsilon)$ where $\epsilon \geq 0.954$, the hierarchical state extension introduces:
\begin{itemize}
    \item Sub-conceptual decomposition function $D: S \rightarrow 2^{\states}$
    \item State lifecycle automaton $L: \states \times \mathcal{C} \rightarrow \states$
    \item DIRAM trace function $T: \states \rightarrow \{0,1\}^{256}$
\end{itemize}
\end{definition}

This extension enables Actors to decompose high-level missions into epistemically validated sub-tasks while maintaining the dimensional innovation property essential to the Actor paradigm.

\section{DIRAM Hardware Fault-Tolerant Architecture}

\subsection{Core State Structure}

The hierarchical state management system anchors to DIRAM's cryptographic memory governance through the following C structure:

\begin{lstlisting}[language=C,caption={DIRAM-backed hierarchical state structure}]
typedef struct {
    uint64_t state_id;
    char parent_state_hash[65];     // SHA-256 trace
    verb_noun_concept_t intent;
    float result_metric;
    float proof_confidence;         // >= 0.954
    state_flag_t status_flag;       // Lifecycle position
    uint8_t error_count;
    uint64_t timestamp;
    diram_state_allocation_t* diram_trace;
} hierarchical_state_t;

typedef enum {
    STATE_TODO = 0x01,
    STATE_DOING = 0x02,
    STATE_DONE = 0x04,
    STATE_BLOCKED = 0x08,
    STATE_ROLLEDBACK = 0x10
} state_flag_t;
\end{lstlisting}

\subsection{Memory Allocation with Trace Linking}

Every state allocation generates a cryptographic receipt ensuring forensic traceability:

\begin{lstlisting}[language=C,caption={DIRAM state allocation implementation}]
diram_state_allocation_t* diram_allocate_state_memory(
    hierarchical_state_t* state,
    const char* intent_tag
) {
    // Enforce epistemic constraint
    if (state->proof_confidence < EPISTEMIC_THRESHOLD) {
        return NULL;
    }
    
    // Generate SHA-256 receipt
    diram_allocation_t* base = diram_alloc_traced(
        sizeof(hierarchical_state_t), intent_tag);
    
    // Link to blockchain for audit trail
    gitraf_blockchain_append_state(
        state->state_id,
        state->parent_state_hash);
    
    return create_state_allocation(base, state);
}
\end{lstlisting}

\section{Task Lifecycle Management with Waterfall Gates}

\subsection{State Transition Automaton}

The lifecycle progression follows a deterministic automaton with epistemic validation at each gate:

\begin{theorem}[Lifecycle Soundness]
For any state $s \in \states$ with confidence $c_s \geq 0.954$, the transition function $L$ guarantees that $L(s, \mathcal{C}) = s'$ implies that the \gls{verifytrace} operation validates the transition $(s \rightarrow s')$ as TRUE.
\end{theorem}

\begin{proof}
Each transition invokes the \gls{audittransition} function which validates the epistemic signature $\Phi$ before permitting state advancement. The DIRAM trace function $T$ generates cryptographic proof of transition validity.
\end{proof}

\subsection{Waterfall Gate Implementation}

\begin{lstlisting}[language=C,caption={Waterfall gate enforcement}]
int enforce_waterfall_gate(
    hierarchical_state_t* state, 
    waterfall_gate_t gate
) {
    switch (gate) {
        case GATE_1_TODO_VALIDATION:
            if (state->proof_confidence < 0.954) {
                state->status_flag = STATE_BLOCKED;
                emit_trace("GATE_1_FAILED", state->state_id);
                return -1;
            }
            break;
            
        case GATE_2_DOING_PROGRESS:
            float ratio = calculate_success_failure_ratio(state);
            if (ratio < 0.5) {  // Below 1:2 threshold
                initiate_cascade_rollback(state);
                return -1;
            }
            break;
            
        case GATE_3_DONE_VERIFICATION:
            emit_verification_proof(state);
            commit_state_to_diram(state);
            break;
    }
    return 0;
}
\end{lstlisting}

\section{Rollback Cascade Protocol}

\subsection{Strategic Rollback Mechanism}

When trial-and-error patterns emerge (error\_count $\geq$ 2), the system initiates the \gls{emitrollback} operation:

\begin{algorithm}
\caption{Cascade Rollback Protocol}
\begin{algorithmic}[1]
\STATE \textbf{Input:} Failed state $s_f$ with confidence $c_f < 0.954$
\STATE \textbf{Output:} Rollback cascade receipt $R$
\STATE $D \leftarrow$ trace-dependency$(s_f)$ \COMMENT{Using \gls{tracedependency}}
\STATE $depth \leftarrow \min(|D|, 5)$ \COMMENT{Limit cascade depth}
\FOR{$d = 0$ \textbf{to} $depth$}
    \STATE $S_d \leftarrow \{s \in D : \text{depth}(s) = d\}$
    \FOR{each $s \in S_d$}
        \STATE $s.confidence \leftarrow s.confidence \times (1 - 0.1d)$
        \STATE $s.status \leftarrow$ STATE\_TODO
        \STATE memoize-delta$(s, c_f)$ \COMMENT{Using \gls{memoizedelta}}
        \STATE generate-receipt$(s)$ \COMMENT{Using \gls{generatereceipt}}
    \ENDFOR
\ENDFOR
\STATE \textbf{return} append-trace$(R)$ \COMMENT{Using \gls{appendtrace}}
\end{algorithmic}
\end{algorithm}

\subsection{Success:Failure Ratio Enforcement}

The system maintains epistemic discipline through continuous ratio monitoring:

\begin{lstlisting}[language=Python,caption={Python implementation of ratio enforcement}]
def assess_state_continuation(self, state):
    """Implements trial-and-improvement with rollback"""
    # Check trial-and-error lock
    if state.confidence < 0.954 and state.error_count >= 2:
        return self._initiate_rollback(state)
    
    # Check success:failure ratio
    ratio = self._calculate_success_ratio(state)
    if ratio < self.rollback_cascade_threshold:  # < 0.5
        return self._strategic_rollback_cascade(state)
    
    # Normal progression
    if state.status_flag == StateFlag.DONE:
        return self._emit_verification_proof(state)
    else:
        return self._update_state(state)
\end{lstlisting}

\section{Actor Sub-ConOps Integration}

\subsection{Alignment with Actor Class Tuple}

The hierarchical state model preserves the Actor's dimensional innovation property while adding structured task management:

\begin{proposition}[Innovation Preservation]
For Actor $\alpha = (S, \mathcal{C}, \Phi, \Psi, \epsilon)$ with hierarchical extension, the dimensional innovation property holds:
\[\exists \tau : S \rightarrow S \text{ where } \tau \notin \text{span}(\mathcal{C}) \implies \exists s \in \states : D(\tau(S)) \ni s\]
\end{proposition}

This ensures that Actor-driven innovations translate to actionable sub-tasks while maintaining epistemic boundaries.

\subsection{Verb-Noun Conceptual Modeling}

Each state intent follows the formalized triplet structure $(V, N, \Phi)$:

\begin{lstlisting}[language=C,caption={Verb-noun concept implementation}]
typedef struct {
    char verb[32];      // Action operation
    char noun[32];      // Domain object  
    float phi_vector[8]; // Epistemic signature
} verb_noun_concept_t;

// Example instantiation
verb_noun_concept_t intent = {
    .verb = "predict",
    .noun = "failure",
    .phi_vector = {0.97, 0.95, 0.98, 0.96, 
                   0.94, 0.99, 0.95, 0.97}
};
\end{lstlisting}

\section{Turing Soundness in Task Decomposition}

\begin{theorem}[Decomposition Completeness]
The hierarchical state system with DIRAM backing achieves Turing-complete task orchestration while maintaining epistemic soundness.
\end{theorem}

\begin{proof}
We construct a correspondence between state transitions and Turing machine computation:
\begin{enumerate}
    \item States in $\states$ encode Turing configurations
    \item Lifecycle transitions simulate state machine evolution
    \item DIRAM provides unbounded memory through linked allocations
    \item Rollback mechanism implements rejection states
    \item The \gls{validateconfidence} operation ensures only sound computations proceed
\end{enumerate}
The 95.4\% threshold prevents non-deterministic branching while cascade protocols enable recovery from computational dead-ends.
\end{proof}

\section{Compliance and Audit Framework}

\subsection{AEGIS-PROOF Traceability}

Every state transition generates auditable proof through:
\begin{itemize}
    \item \gls{commitstate}: Persistence with cryptographic receipt
    \item \gls{anchorhardware}: Physical memory binding for forensics
    \item \gls{computeratio}: Continuous success metric validation
\end{itemize}

\subsection{NASA-STD-8739.8 Adherence}

The system satisfies safety-critical requirements through:
\begin{enumerate}
    \item \textbf{Deterministic Execution}: State transitions follow formal automaton
    \item \textbf{Bounded Resources}: DIRAM enforces $\epsilon(x) \leq 0.6$ constraint
    \item \textbf{Graceful Degradation}: Cascade rollback prevents catastrophic failure
    \item \textbf{Formal Verification}: All paths traceable through SHA-256 receipts
\end{enumerate}

\section{Production Deployment Architecture}

\begin{lstlisting}[language=Python,caption={Complete orchestrator implementation}]
class ActorSubConOpsOrchestrator:
    """Production-ready hierarchical task orchestration"""
    
    def __init__(self):
        self.epistemic_threshold = 0.954
        self.rollback_cascade_threshold = 0.5
        self.diram = DIRAMInterface()
        
    def process_mission(self, actor, mission):
        # Decompose using dimensional innovation
        states = self.decompose_mission(actor, mission)
        
        # Process each state through lifecycle
        for state in states:
            while state.status_flag != STATE_DONE:
                transition = self.process_state_lifecycle(state)
                
                if transition == StateTransition.ROLLBACK:
                    self.handle_cascade_recovery(state)
                elif transition == StateTransition.BLOCKED:
                    self.resolve_dependencies(state)
                    
        return self.compile_mission_proof(states)
\end{lstlisting}

\section{Conclusion}

The hierarchical Actor-orchestrated state management system represents deployed infrastructure achieving self-correcting AI orchestration through:
\begin{itemize}
    \item DIRAM-backed memory governance with cryptographic traceability
    \item 95.4\% epistemic validation threshold enforcement
    \item Strategic rollback cascades maintaining 1:2 success:failure ratios
    \item Verb-noun conceptual modeling for semantic task representation
    \item Waterfall gate compliance for systematic validation
\end{itemize}

This architecture operates continuously across OBINexus deployments, transforming Actor-level dimensional innovations into tractable, verifiable sub-tasks while maintaining the mathematical rigor demanded by safety-critical AI systems.

\printglossary[title={Verb-Noun Concept Glossary},toctitle={Glossary}]

\end{document}